%%%%%%%%%%%%%%%%%%%%%
% Documento maestro %
%%%%%%%%%%%%%%%%%%%%%
\documentclass{fime}

%%%%%%%%%%%%%%%%%%%%%%%%%%%%%%%%%%%%%%%%%%%
% Cargando paquetes y definiendo opciones %
%%%%%%%%%%%%%%%%%%%%%%%%%%%%%%%%%%%%%%%%%%%
% Aquí puedes cargar los paquetes que vas a usar. La clase
% fime ya incluye babel, inputenc, graphicx y los de la AMS.
% Cargar un paquete está a tu libertad (y responsabilidad).
\usepackage{hyperref}
    \hypersetup{breaklinks=true,colorlinks=true,linkcolor=black,citecolor=black,urlcolor=black}

%%%%%%%%%%%%%%%%%%%%%
% Definiendo campos %
%%%%%%%%%%%%%%%%%%%%%
\def\titulo{Título de la tesis}
\def\autor{Fulano D'Tal}
\def\matricula{XXXXXX}
\def\grado{Maestría en Logística y Cadena de Suministro}
% En caso de que el grado tenga orientación o especialidad llenar el siguiente
% campo dejando un ESPACIO INICIAL, en caso contrario, dejar vacío
\def\orientacion{ con orientación en Diseño y Análisis}
% Coloca el mes con mayúscula inicial
\def\fecha{Octubre 2013}

\def\asesor{Dr. Asesor}
\def\revisorA{Nombre del revisor A}
\def\revisorB{Nombre del revisor B}
% En el caso de que tu tesis sea de doctorado activa la variable cambiándola a \doctoradotrue
% y define tus otros dos revisores
\newif\ifdoctorado\doctoradofalse
\def\revisorC{Nombre del revisor C}
\def\revisorD{Nombre del revisor D}
% El visto bueno siempre va
\def\vobo{Dr. Simón Martínez Martínez}

%%%%%%%%%%%%%%%%%%%%%%%
% Inicia el documento %
%%%%%%%%%%%%%%%%%%%%%%%
\begin{document}

\frontmatter
\pagestyle{main}

%%%%%%%%%%%%%%%%%%%%%%%%
% Primer portada: UANL %
%%%%%%%%%%%%%%%%%%%%%%%%
\thispagestyle{empty}

\begin{scshape}
\begin{center}
	{\Large \uanl} \\[5mm]
	{\large \fime} \\[5mm]
	{\large \depg}
	\vskip15mm
	\includegraphics[height=55mm]{uanl}
	\vskip12mm
	\begin{tabular}{p{11cm}}
		\centering
		{\large \titulo}
	\end{tabular}
	\vskip7mm
	{por}\\[7mm]
	{\large \autor}\\[7mm]
        % Si en tu posgrado la tesis no es opcional (como sí lo es en licenciatura), no modifiques esta línea:
	{como requisito parcial para obtener el grado de}\\[3mm]
	\MakeUppercase{\grado}\\
	\orientacion
	\vfill
	\fecha
\end{center}
\end{scshape}

%%%%%%%%%%%%%%%%%%%%%%%%%
% Segunda portada: FIME %
%%%%%%%%%%%%%%%%%%%%%%%%%
\newpage
\thispagestyle{empty}

\begin{scshape}
\begin{center}
	{\Large\uanl} \\[5mm]
	{\large\fime} \\[5mm]
	{\large\depg}
	\vskip16mm
	\includegraphics[height=50mm]{fime}
	\vskip16mm
	\begin{tabular}{p{11cm}}
		\centering
		{\large \titulo}
	\end{tabular}
	\vskip7mm
	{por}\\[7mm]
	{\large \autor}\\[7mm]
        % Si en tu posgrado la tesis no es opcional (como sí lo es en licenciatura), no modifiques esta línea:
	{como requisito parcial para obtener el grado de}\\[3mm]
	\MakeUppercase{\grado}\\
	\orientacion
	\vfill
	\fecha
\end{center}
\end{scshape}

%%%%%%%%%%%%%%%%%%%%%%%%%%%%%
% Carta del comité de tesis %
%%%%%%%%%%%%%%%%%%%%%%%%%%%%%
\newpage
\thispagestyle{empty}
\enlargethispage{8mm}

{\renewcommand{\baselinestretch}{1.1}\selectfont
\begin{center}\vspace*{-25mm}\hspace*{-10mm}
\begin{minipage}{170.5mm}
\hspace{-1.5mm}\includegraphics[height=20mm]{uanl}\hfill\raise1mm\hbox{\includegraphics[height=18.5mm]{fime}}
\hrule\vspace{0.5mm}
\scalebox{.5}{\MakeUppercase{\uanl}}\hfill\scalebox{.5}{\MakeUppercase{\fime}}\medskip
\end{minipage}
\vskip4mm{\sc\large\uanl\\\fime\\[3pt]\depg}\vskip6mm
\end{center}

Los miembros del Comité de Tesis recomendamos que la Tesis <<\titulo>>, realizada por el alumno \autor, con número de matrícula \matricula, sea aceptada para su defensa como requisito para obtener el grado de \grado\orientacion.
\ifdoctorado\vskip10mm\else\vskip8mm\fi

\begin{center}
El Comité de Tesis\\
\ifdoctorado\vskip12mm\else\vskip25mm\fi
\ifdoctorado{%%%
\begin{tabular}{p{37mm}p{21mm}p{12mm}p{21mm}p{37mm}}
	\cline{2-4}
	& \multicolumn{3}{c}{\asesor} & \\
	& \multicolumn{3}{c}{Asesor}  & \\[15mm]
	\cline{1-2} \cline{4-5}
	\multicolumn{2}{c}{\revisorA} & & \multicolumn{2}{c}{\revisorB} \\
	\multicolumn{2}{c}{Revisor}   & & \multicolumn{2}{c}{Revisor}   \\[17mm]
	\cline{1-2} \cline{4-5}
	\multicolumn{2}{c}{\revisorC} & & \multicolumn{2}{c}{\revisorD} \\
	\multicolumn{2}{c}{Revisor}   & & \multicolumn{2}{c}{Revisor}   \\[2mm]
	& \multicolumn{3}{c}{Vo. Bo.} & \\[14mm]
	\cline{2-4}
	& \multicolumn{3}{c}{\vobo}   & \\
	& \multicolumn{3}{c}{Subdirector de Estudios de Posgrado}   & \\ &&&&
\end{tabular}
}\else{%%%
\begin{tabular}{p{37mm}p{21mm}p{12mm}p{21mm}p{37mm}}
	\cline{2-4}
	& \multicolumn{3}{c}{\asesor} & \\
	& \multicolumn{3}{c}{Asesor}  & \\[19mm]
	\cline{1-2} \cline{4-5}
	\multicolumn{2}{c}{\revisorA} & & \multicolumn{2}{c}{\revisorB} \\
	\multicolumn{2}{c}{Revisor}   & & \multicolumn{2}{c}{Revisor}   \\[2mm]
	& \multicolumn{3}{c}{Vo. Bo.} & \\[17mm]
	\cline{2-4}
	& \multicolumn{3}{c}{\vobo}   & \\
	& \multicolumn{3}{c}{Subdirector de Estudios de Posgrado}   & \\ &&&&
\end{tabular}
}\fi%%%
\vfill
\snnl, \MakeLowercase{\fecha}
\end{center}
\vfill
\begin{flushright}
\begin{minipage}{68mm}
\sffamily\tiny\color{gray}
Ciudad Universitaria,\hfill Pedro de Alba s/n,\hfill C.P. 66455,\hfill A.P. 076 Suc.\! ``F''\\
San Nicolás de los Garza,\hfill Nuevo León,\hfill México.\hfill Tels.: 81 8332 0903\\
Conm.: 81 8329 4020.\; Fax: 81 8332 0904.
\end{minipage}
\end{flushright}
}

% Dedicatoria

\thispagestyle{empty}
\vspace*{17mm}

\begin{flushright}
\begin{itshape}

Aquí puedes poner tu dedicatoria\\
si es que tienes una.\bigskip\bigskip

Si no tienes una, puedes borrar\\
la línea \verb+% Dedicatoria

\thispagestyle{empty}
\vspace*{17mm}

\begin{flushright}
\begin{itshape}

Aquí puedes poner tu dedicatoria\\
si es que tienes una.\bigskip\bigskip

Si no tienes una, puedes borrar\\
la línea \verb+% Dedicatoria

\thispagestyle{empty}
\vspace*{17mm}

\begin{flushright}
\begin{itshape}

Aquí puedes poner tu dedicatoria\\
si es que tienes una.\bigskip\bigskip

Si no tienes una, puedes borrar\\
la línea \verb+\include{Dedicatoria}+ en el\\
archivo \texttt{MiTesis.tex} pues no es obligatoria.

\end{itshape}
\end{flushright}

+ en el\\
archivo \texttt{MiTesis.tex} pues no es obligatoria.

\end{itshape}
\end{flushright}

+ en el\\
archivo \texttt{MiTesis.tex} pues no es obligatoria.

\end{itshape}
\end{flushright}



\tableofcontents
\listoffigures
\listoftables

%Agradecimientos

\chapter{Agradecimientos}
\markboth{Agradecimientos}{}

Aquí puedes poner tus agradecimientos. (No olvides agradecer a tu comité de tesis, a tus profesores, a la facultad y a CONACyT en caso de que hallas sido beneficiado con una beca).

%Resumen

\chapter{Resumen}
\markboth{Resumen}{}

{\renewcommand{\baselinestretch}{1.1}\selectfont
{\setlength{\leftskip}{10mm}
\setlength{\parindent}{-10mm}

\autor.

Candidato para obtener el grado de \grado\orientacion.

\uanl.

\fime.

Título del estudio: \textsc{\titulo}.

\noindent Número de páginas: \pageref*{lastpage}.}

%%% Comienza a llenar aquí
\paragraph{Objetivos y método de estudio:}
Aquí debes poner tus objetivos y métodos de estudio. (Este es el formato).

\paragraph{Contribuciones y conlusiones:}
Y aquí tus contribuciones y conclusiones. (También es parte del formato).

\bigskip\noindent\begin{tabular}{lc}
\vspace*{-2mm}\hspace*{-2mm}Firma del asesor: & \\
\cline{2-2} & \hspace*{1em}\asesor\hspace*{1em}
\end{tabular}}



\mainmatter
\pagestyle{fime}

%%% Haz un documento para cada capítulo
\chapter{Este es un capítulo}

\section{Esta es una sección}

Quiero escribir un texto largo para ver la sangría que ha sido colocada en el formato del párrafo. También ha sido declarado el espacio entre párrafos, pero puedes agregar espacio adicional cuando se requiera (entre algunos objetos o entornos quizá llegue a ser necesario).

Y aquí quiero escribir otro texto largo para ver el espacio entre párrafos colocado por defecto.

\subsection{Entornos}

Ahora mostraremos cómo quedan algunos entornos:

\paragraph{Una tabla:} El estilo de tesis de la UANL requiere que la descripción de las tablas sea de tipo encabezado, así que recuerda colocar el {\tt caption} antes que la tabla.

\begin{table}[ht]
	\centering
	\caption{Una tablita}
	\label{tab:una-tablita}
		\begin{tabular}{|c|c|}
			\hline
			a & b \\
			\hline
			c & d \\
			\hline
		\end{tabular}
\end{table}

\paragraph{Una figura:} Ya ha sido cargado el paquete {\tt graphicx} en {\tt fime.cls}, así que podemos incluir gráficas. Coloca todas tus figuras en la carpeta {\tt Figuras} ya que esa es la ubicación en que las buscará {\tt fime.cls}, además, de esa manera tendrás un poco más de orden en tus archivos. El estilo de tesis de la UANL requiere que la descripción de las figuras sea al calce, así que recuerda colocar el {\tt caption} después que la figura.

\begin{figure}[htp]
	\centering
		\includegraphics[width=4cm]{uanl.eps}
	\caption{El escudo de la UANL}
	\label{fig:uanl}
\end{figure}

\subsection{Matemáticas}

Ahora mostraremos cómo quedan algunos objetos matemáticos: una ecuación y una definición.

La ecuación \eqref{eq:obvia} es obvia.
\begin{equation}
	x + x = 2x
	\label{eq:obvia}
\end{equation}

Ahora una definición. En {\tt fime.cls} se definen otros entornos como teorema, lema, proposición, ejemplo, etc. Echa un vistazo y si falta el que necesitas, ¡agrégalo!

\begin{definicion}
	Un número entero es primo si y sólo si tiene exactamente cuatro divisores distintos.
\end{definicion}

Agreguemos otro entorno para que se vea la numeración de estos.

\begin{teorema}[Fermat]
	No existen $x,y,z \in \mathbb{Z}$ (que no sean los triviales) tales que $x^n + y^n = z^n$ para $n \geq 3$.
\end{teorema}


\appendix
%%% Haz un documento para cada apéndice, si es que tienes
\chapter{Este es un apéndice}

\section{Citas bibliográficas}

En princicpio tienes total libertad de incluir tu bibliografía con el entorno {\tt thebibliography} nativo de \LaTeX{} o mediante la herramienta \textsc{Bib}\TeX. En caso de que optes por esto último (recomendado), puedes usar alguno de los archivos {\tt mighelbib.bst} o {\tt mighelnat.bst} incluidos en el paquete {\tt Tesis-FIME}, pues sus diseños están basados en el estilo bibliográfico estándar del español, además de que armoniza con el estilo de tesis provisto por {\tt fime.cls}.

El estilo bibliográfico {\tt mighelbib} es numérico, es decir cita con un número entre corchetes, por ejemplo una cita \verb+\cite{Dan82}+ genera una etiqueta del tipo [13], mientras que el estilo {\tt mighelnat} es tipo autor-año y requiere que el paquete {\tt natbib} sea cargado (sin opciones) para su correcto funcionamiento, cita con el apellido del autor y el año, por ejemplo una cita \verb+\citet{Dan82}+ genera una etiqueta del tipo Dantzig (1982), mientras que una cita \verb+\citep{Dan82}+ genera una etiqueta del tipo (Dantzig, 1982).

Como muestra del estilo, unas citas: un libro clásico de programación lineal \citep{Baz04} y un documento histórico \citep{Dan82}. Para saber un poco más del uso de \textsc{Bib}\TeX, se puede leer \citep{Mat11}.

\section{Comillas}

El objetivo de esta sección era provocar otra página para que se vea el encabezado. Pero aprovechamos para decir que la clase {\tt fime.cls} carga el paquete {\tt babel} con la opción {\tt spanish}, por lo que cambiará automáticamente los dobles signos $<<$ y $>>$ por << y >>. Estas comillas angulares son las correctas en el idioma español, y son las que se usan en la clase {\tt fime.cls}, por lo que se sugiere sean las usadas en el texto cada que quieras <<entrecomillar>> algo.



\backmatter
\pagestyle{main}

% Aquí va la bibliografía, puedes usar el entorno de LaTeX (thebibliography)
% o la herramienta BibTeX. En caso de que optes por BibTeX, puedes usar
% alguno de los archivos de estilo (mighelbib.bst o mighelnat.bst) incluidos
% en el paquete, cuyos diseños armonizan con el diseño de tesis provisto por
% fime.cls. Para muestra, basta un botón:
\bibliographystyle{mighelnat}
\bibliography{MiBiblio}

\label{lastpage}
%Autobiografia

\chapter*{Resumen autobiográfico}
\thispagestyle{empty}

\begin{center}
\autor

Candidato para obtener el grado de\\
\grado\\
\orientacion\bigskip

\uanl\\
\fime\bigskip

Tesis:\\
\textsc{\large\titulo}
\end{center}\bigskip

%Aquí va tu historia
Aquí va tu historia. Recuerda que debe incluir: lugar y fecha de nacimiento, nombre de los padres, escuelas y universidades en las que se graduó después de la preparatoria, títulos o grados obtenidos (no incluir los estudios que se están concluyendo), experiencia profesional y organizaciones profesionales a las que pertenece (no incluir lista de publicaciones).


\end{document}

